\subsection{Joint prediction for sequences, grouped units, and multivariate responses}
\label{sec:prediction-extensions}
The fixed-segment predictive formula in Proposition~\ref{prop:predictive-paper} supports several input structures, but \emph{joint} scoring must preserve dependence induced by a shared segment parameter. A product of separately integrated single-observation predictive densities is generally not the integrated density of their joint collection.

\begin{proposition}[Joint predictive integration within a fitted segment]
\label{prop:joint-predictive-paper}
Condition on a fitted partition $t$ and conjugate posterior $p(\theta_B\mid y_B)$ for each of its segments $B$. If new observations $D_B^{\rm new}$ assigned to $B$ are conditionally independent given $\theta_B$ and have sufficient summaries $(S_B^{\rm new},W_B^{\rm new},H_B^{\rm new})$, then
\begin{equation}
p(D_B^{\rm new}\mid y_B,t)
=e^{H_B^{\rm new}}\frac{Z(\alpha_B+S_B^{\rm new},\beta_B+W_B^{\rm new})}{Z(\alpha_B,\beta_B)}.
\label{eq:joint-predictive-block}
\end{equation}
The full fixed-partition joint predictive density is the product of these contributions over \emph{distinct segments}, not in general over individual observations within one segment.
\end{proposition}
\begin{proof}
\textbf{Step 1 (conditional likelihood).} Given $\theta_B$, independence gives $p(D_B^{\rm new}\mid\theta_B)=\prod_{r\in D_B^{\rm new}}p(y_r\mid\theta_B,d_r)$.
\textbf{Step 2 (aggregate once).} The sufficient-statistic representation collects all new observations into one exponential-family kernel with summaries $S_B^{\rm new},W_B^{\rm new},H_B^{\rm new}$.
\textbf{Step 3 (integrate once).} Multiply this kernel by the training posterior and integrate $\theta_B$, yielding Equation~\eqref{eq:joint-predictive-block}. It differs from integrating $\theta_B$ separately for each observation, which would spuriously treat new observations as having independent latent parameters.
\textbf{Step 4 (between-segment factorization).} Under the model assumption of conditionally independent segment parameters, the posterior factorizes across distinct segments given $t$. Therefore the joint predictive density factors over the distinct segments, completing the claim.
\end{proof}


\begin{corollary}[Closed-form predictive distributions in three conjugate families]
\label{cor:predictive-families-paper}
Conditional on a fitted segment and its conjugate posterior, the following are ordinary posterior-predictive sampling distributions:
\begin{enumerate}
\item For Gaussian $\mu\mid y_B\sim\mathcal N(\nu_B,V_B)$ and new observations with known variances $\sigma^2/w_r$, a vector of $m$ new observations in the same segment has joint law
\[
Y^{\rm new}\mid y_B,t\sim\mathcal N_m\!\left(\nu_B\mathbf1_m,\;V_B\mathbf1_m\mathbf1_m^\top+\operatorname{diag}(\sigma^2/w_1,\ldots,\sigma^2/w_m)\right).
\]
\item For Poisson $\lambda\mid y_B\sim\operatorname{Gamma}(a_B,b_B)$ and a new count at exposure $e>0$, the predictive law is negative Binomial with
\[
p(y^{\rm new}=z\mid y_B,e)=\frac{\Gamma(a_B+z)}{\Gamma(a_B)z!}\left(\frac{b_B}{b_B+e}\right)^{a_B}\left(\frac{e}{b_B+e}\right)^z.
\]
\item For Binomial $p\mid y_B\sim\operatorname{Beta}(a_B,b_B)$ and $m$ future trials, the predictive count law is Beta--Binomial with
\[
p(y^{\rm new}=z\mid y_B,m)=\binom mz\frac{B(a_B+z,b_B+m-z)}{B(a_B,b_B)}.
\]
\end{enumerate}
\end{corollary}
\begin{proof}
\textbf{Step 1 (Gaussian joint law).}
Write each new observation as $Y_r=\mu+\epsilon_r$ with independent $\epsilon_r\sim\mathcal N(0,\sigma^2/w_r)$ conditional on $\mu$, and with $\mu\sim\mathcal N(\nu_B,V_B)$ from the training posterior. Gaussian addition gives mean $\nu_B\mathbf1_m$ and covariance $V_B\mathbf1_m\mathbf1_m^\top+\operatorname{diag}(\sigma^2/w_r)$. The rank-one term shows why multiplying separate Gaussian marginals would be incorrect for joint scoring.
\textbf{Step 2 (Poisson--Gamma mixture).}
Integrate $e^{-e\lambda}(e\lambda)^z/z!$ against the Gamma posterior density. The Gamma integral has shape $a_B+z$ and rate $b_B+e$, which simplifies to the displayed negative-Binomial probability mass function.
\textbf{Step 3 (Beta--Binomial mixture).}
Integrate $\binom mzp^z(1-p)^{m-z}$ against $\operatorname{Beta}(a_B,b_B)$. Combining exponents produces $B(a_B+z,b_B+m-z)/B(a_B,b_B)$, establishing the third expression.
\end{proof}

\paragraph{Pointwise sequences (Case A).}
New pointwise observations carry coordinates, outcomes, and the family-specific descriptors. Under an exported partition, coordinates select segments and the data are aggregated \emph{within} those segments before applying Equation~\eqref{eq:joint-predictive-block}. An out-of-range coordinate requires a declared extrapolation model; a coordinate gap is not an exposure by default.

\paragraph{Set-valued units (Case B).}
A unit may contain multiple pairs $(x_{ur},y_{ur})$ on a physical support interval. Its segment-specific internal observations can be scored jointly using the same sufficient-statistic aggregation. If several units in a single dataset share segment parameters, the dataset score aggregates \emph{across all such units before integrating}. Multiplying separately marginalized unit likelihoods is justified only under an additional model in which the units have independent latent segment parameters; it is otherwise a composite score.

\paragraph{Vector-valued responses (Case C).}
When response coordinates are conditionally independent given distinct component-specific segment parameters, the joint segment evidence and predictive density factor across response dimensions. For a genuinely multivariate exponential family with shared or correlated parameters, the joint sufficient statistics and conjugate normalizer must be used instead. The global ordered-partition recursion depends only on the resulting joint segment score.

\paragraph{Group membership and segmentation averaging.}
If $M_g$ defines a \emph{normalized} predictive data density and $\pi_g$ is a prior group probability, then $p(g\mid D^{\rm new})\propto\pi_gp(D^{\rm new}\mid M_g)$ is a Bayesian group posterior. Scores from the finite latent-group MM criterion in Equation~\eqref{eq:latent-group-objective}, which are not asserted to be normalized data densities, do not automatically admit this interpretation. Likewise, true segmentation-averaged prediction sums joint predictive contributions over $p(k,t\mid y)$; prediction conditional on a chosen joint-MAP partition is a different object.

\paragraph{Calibration.}
For discrete outcomes, randomize the probability integral transform within the probability mass at the observed outcome. For continuous outcomes, PIT values are marginally uniform under a calibrated predictive law, but they are generally dependent within a sequence. Report sequence-level calibration or simulation-based reference distributions rather than applying an independent-uniform goodness-of-fit reference blindly. Held-out predictive scores need not rank candidate $k$ in the same order as training posterior count probabilities.
